% !TEX program = xelatex
% ===========================================================================
%  : HallucinationCompactness
%  The Auditability Principle:
%  Everything Real is Auditable — What Cannot Be Audited Is Either
%  Illusion or Compactness-Inseparable
% ===========================================================================
%  SCX Capstone Philosophy Paper
%  The Auditability Principle — the SCX analog of the Church-Turing Thesis
% ===========================================================================
\documentclass[12pt,a4paper]{article}
% ===========================================================================
%  Packages
% ===========================================================================
\usepackage{amsmath,amssymb,amsthm}
\usepackage{mathtools}
\usepackage{mathrsfs}
\usepackage{bbm}
\usepackage{bm}
\usepackage{graphicx}
\usepackage{xcolor}
\usepackage{hyperref}
\usepackage{enumitem}
\usepackage{booktabs}
\usepackage{tikz}
\usepackage{tikz-cd}
\usepackage{float}
\usepackage{geometry}
\usepackage{cleveref}
\usepackage{physics}
\geometry{margin=2.5cm,top=2cm,bottom=2cm}
% ===========================================================================
%  Hyperref setup
% ===========================================================================
\hypersetup{
    colorlinks=true,
    linkcolor=blue,
    citecolor=blue,
    urlcolor=blue,
    pdftitle={:},
    pdfauthor={SCX Capstone Philosophy},
    pdfsubject={The Auditability Principle — SCX Analog of Church-Turing Thesis},
% ===========================================================================
%  Theorem environments
% ===========================================================================
\theoremstyle{plain}
\newtheorem{theorem}Theorem[section]
\newtheorem{lemma}[theorem]Lemma
\newtheorem{proposition}[theorem]{Proposition}
\newtheorem{corollary}[theorem]Corollary
\newtheorem{conjecture}[theorem]Conjecture
\theoremstyle{definition}
\newtheorem{definition}[theorem]Definition
\newtheorem{example}[theorem]Example
\newtheorem{remark}[theorem]Remark
\newtheorem{principle}[theorem]Principle
\theoremstyle{remark}
\newtheorem*{note}Note
% ===========================================================================
%  Custom commands
% ===========================================================================
\newcommand{\R}{\mathbb{R}}
\newcommand{\C}{\mathbb{C}}
\newcommand{\N}{\mathbb{N}}
\newcommand{\Z}{\mathbb{Z}}
\newcommand{\Q}{\mathbb{Q}}
\newcommand{\E}{\mathbb{E}}
\newcommand{\Pbb}{\mathbb{P}}
\newcommand{\Var}{\mathrm{Var}}
\newcommand{\D}{\mathcal{D}}
\newcommand{\G}{\mathcal{G}}
\newcommand{\T}{\mathcal{T}}
\newcommand{\M}{\mathcal{M}}
\newcommand{\I}{\mathcal{I}}
\newcommand{\F}{\mathcal{F}}
\newcommand{\A}{\mathcal{A}}
\newcommand{\Ccal}{\mathcal{C}}
\newcommand{\SCX}{\textsf{SCX}}
\newcommand{\g}{\mathbf{g}}
\newcommand{\x}{\mathbf{x}}
\newcommand{\X}{\mathcal{X}}
\newcommand{\Obs}{\mathcal{O}}
\newcommand{\ClaimSpace}{\mathfrak{C}}
\newcommand{\AuditBoundary}{\partial\mathfrak{A}}
\newcommand{\Null}{\varnothing}
\DeclareMathOperator{\argmax}{arg\,max}
\DeclareMathOperator{\argmin}{arg\,min}
\DeclareMathOperator{\sgn}{sgn}
\DeclareMathOperator{\diag}{diag}
\DeclareMathOperator{\Yajie}{Yajie}
\DeclareMathOperator{\Cercis}{Cercis}
\DeclareMathOperator{\Situs}{Situs}
\DeclareMathOperator{\Audit}{Audit}
\DeclareMathOperator{\ObsSet}{Obs}
\DeclareMathOperator{\Claim}{Claim}
\DeclareMathOperator{\Ver}{Ver}
\DeclareMathOperator{\Comp}{Comp}
\DeclareMathOperator{\supp}{supp}
% Honest critique
\newcommand{\honestcrit}[1]{%
    \textcolor{red}{\textbf{[Honest Strike]#1}}%
% Proof-status markers
\newcommand{\rigorFull}{\textnormal{\textbf{[Rigorous Proof]}}}
\newcommand{\rigorPartial}{\textnormal{\textbf{[Partial Proof]}}}
\newcommand{\rigorSketch}{\textnormal{\textbf{[Proof Sketch]}}}
\newcommand{\openproblem}{\textnormal{\textbf{[Open Problem]}}}
% Special environments for AP
\newcommand{\AP}{\mathsf{AP}}
\newcommand{\CI}{\mathsf{CI}}
\newcommand{\UD}{\mathsf{UNDECLARED}}
\newcommand{\NOISY}{\mathsf{NOISY}}
% ===========================================================================
%  Title
% ===========================================================================
\title{
    \textbf{: \\[4pt]}\\[12pt]
    {\large The Auditability Principle:\\
    Everything Real is Auditable — What Cannot Be Audited\\
    Is Either Illusion or Compactness-Inseparable}\\[8pt]
    {\normalsize SCX Capstone Philosophy Paper}\\[4pt]
    {\footnotesize The SCX Analog of the Church-Turing Thesis}
\author{
    SCX Capstone Research Group\\[2pt]
    \texttt{Xiaogan Supercomputing Center}\\[4pt]
    \small\textit{Version 1.0 — July 2026}
\date{}
% ===========================================================================
\begin{document}
\maketitle
% ===========================================================================
%  Abstract
% ===========================================================================
\begin{abstract}
\noindent
\textbf{Abstract: } SCX\\textbf{Auditability Principle}（Auditability Principle, AP）——SCXChurch-Turing: 
\textit{" $g=0$}$ ."}
Proof——, .Proof, : (1) \textbf{Hallucination}——, YajieNOISY, ; (2) \textbf{Compactness}—— $M\to\infty$  $t\to\infty$ , awaiting, （G\"odel）.AP, AP.
Compactness: Compactness.\textbf{}.SCX——\textit{}., : .
\vspace{8pt}
\noindent
\textbf{English Abstract:} We propose the \\textbf{Auditability Principle} (AP) — the SCX analog of the Church-Turing Thesis:
\textit{"Every genuine claim can be expressed as a statement about $g=0$ for some observable system, and this statement can be verified by $M>1$ independent observers with bounded error $e^{-2M\Delta^2}$."}
The principle is not provable — it is a thesis, serving as a taxonomic principle to characterize the boundaries of auditing. We demonstrate that all allegedly unauditable claims fall into exactly two categories: (1) \textbf{Illusion} — claims with no observable consequences, classified by Yajie as NOISY, never having been real claims at all; (2) \textbf{Compactness-Inseparable} — claims whose verification requires $M\to\infty$ or $t\to\infty$, operationally equivalent to unauditable yet logically well-formed (analogous to G\"odel sentences). These are not counterexamples to AP — they are boundary conditions that AP as a taxonomic framework correctly predicts.
We establish the compactness boundary of audit theory: the frontier between the space of auditable claims and the space of unauditable claims is determined by the compactness properties of the claim space. This boundary is itself \textbf{approximately detectable}. SCX cannot audit everything — it can audit everything \textit{that can be audited}.
\vspace{8pt}
\noindent
\textbf{Keywords:} Auditability Principle,  SCX framework,  Church-Turing thesis,  compactness boundary,  audit-undecidability,  Yajie classification,  G\"odel sentences,  operational equivalence,  illusion,  compactness-inseparable,  complete theory
\noindent
\textbf{Keywords:
\end{abstract}
\newpage
\tableofcontents
\newpage
% ===========================================================================
%  PREAMBLE — WHY THIS PAPER EXISTS
% ===========================================================================
\section{:}
\section{Preamble: Why the Auditability Principle Is Necessary}
\subsection{SCXChurch-Turing}
\subsection{SCX Needs Its Church-Turing Moment}
\textit{SCX？}
\textit{SCX\\textbf{}？, ？, ？}
\vspace{6pt}
\noindent
\textbf{English:} The SCX framework has developed a formidable mathematical skeleton: potential surface theory, Yajie consensus auditing, the Situs manifold, compactness analysis, instanton detection, singularity theory, and more. Audit techniques grow ever more refined — from pointwise Yajie auditing to persistent homology, from gauge theory to quantum frameworks — but we have never answered the most fundamental question:
\textit{What are the limits of SCX auditing?}
Or, more pointedly:
\textit{Is there anything SCX cannot audit \textbf{in principle}? If so, what? If not, why?}
\vspace{6pt}
2030.Turing, Church, G\"odel, Post, Kleeneawaiting, ChurchTuring: \textit{Turing.}Proof——""——."Turing", ProofawaitingTuring, .
\vspace{4pt}
\noindent
\textit{Computation theory faced an almost identical problem in the 1930s. Turing, Church, G\"odel, Post, Kleene, and others each proposed formal definitions of computability, and then Church and Turing independently proposed the famous thesis: every intuitively computable function is Turing-computable. This thesis is not provable — "intuitively computable" is not a formal concept — but it has withstood nearly a century of testing. Every attempt to define a "stronger-than-Turing" model of computation has either been proven equivalent to Turing machines, or falsified.}
\vspace{6pt}
\vspace{4pt}
\noindent
\textit{SCX needs its own "Church-Turing moment." This paper provides that moment.}
\subsection{Auditability Principle}
\subsection{Status of the Auditability Principle}
AP\textbf{}（thesis）.: 
\begin{enumerate}[label=(\roman*)]
    \item SCX\textbf{}——AP, ; 
    \item \textbf{}——"", AP: HallucinationCompactness？
    \item \textbf{}——, .
\end{enumerate}
\vspace{4pt}
\noindent
\textit{The Auditability Principle (AP) is not a theorem. It has no proof — not because a proof has not been found, but because its premise ("genuine claim") is not itself a formal concept, just as the premise of the Church-Turing thesis ("intuitively computable") is not a formal concept.}
\textit{AP is a \textbf{thesis}. Its role is to: (i) delineate the \textbf{applicability boundary} of the SCX audit framework — telling us what kinds of claims fall within audit scope; (ii) provide a \textbf{diagnostic tool} — if a claim is said to be "unauditable," AP tells us what to check: illusion or compactness-inseparable? (iii) assert \textbf{theoretical completeness} — a theory that knows its own limits is more complete than one that claims omnipotence.}
\subsection{}
\subsection{Structure of This Paper}
\textbf{}（tightening spiral）——Core: 
\begin{enumerate}[label=\S\arabic*:]
    \item \textbf{Auditability Principle（AP）}——, , Church-Turing; 
    \item \textbf{""（Why）}——HallucinationCompactness; 
    \item \textbf{Compactness}——CompactnessVerdict; 
    \item \textbf{awaiting}——$M\to\infty$ awaiting; 
    \item \textbf{}——SCX, ; 
    \item \textbf{Honest Strike}——SCX, .
\end{enumerate}
\vspace{4pt}
\noindent
\textit{This paper is structured as a \textbf{tightening spiral} — each section deepens the same core insight on the foundation laid by the previous section: \S1: The AP — formal statement, formalization, correspondence with Church-Turing; \S2: Two "exceptions" — complete taxonomy of illusion and compactness-inseparable; \S3: The compactness boundary — rigorous correspondence from first-order logic compactness to audit-undecidability; \S4: Operational equivalence — $M\to\infty$ is operationally equivalent to unauditable; \S5: What this means — implications for SCX practice, philosophy, and future; \S6: The final honest hit — SCX knows what it can do, and what it cannot.}
\newpage
% ===========================================================================
%  SECTION 1: THE AUDITABILITY PRINCIPLE
% ===========================================================================
\section{Auditability Principle}
\section{The Auditability Principle}
\subsection{}
\subsection{Formal Statement}
\begin{principle}[Auditability Principle — The Auditability Principle]\label{pr:AP}
 $\X$ , $\ClaimSpace$ （claim）.: 
\begin{enumerate}[label=(\roman*)]
    \item \textbf{ (Expressibility):}  $c \in \ClaimSpace$  $S \in \X$  $g_S$ : 
    \begin{equation}\label{eq:AP_express}
        c \iff (g_S = 0) \quad \text{} \quad S \in \X.
    \end{equation}
    \item \textbf{ (Verifiability):}  $M > 1$ Observers: 
    \begin{equation}\label{eq:AP_verify}
        \Pbb(\text{ $M$}) \leq e^{-2M\Delta^2},
    \end{equation}
     $\Delta = \min_m |\E[\hat{g}_m] - g_S|$ .
    \item \textbf{ (Boundary Condition):}  $c$  $M$  $t$ ,  $c$ \\
    (a) \textbf{Hallucination}（, $\ObsSet(c) = \Null$）, \\
    (b) \textbf{Compactness}（ $M \to \infty$  $t \to \infty$）.
\end{enumerate}
\end{principle}
\vspace{6pt}
\noindent
\textbf{English:} Let $\X$ be the set of all observable systems and $\ClaimSpace$ the set of all claims. Then: (i) \textit{Expressibility:} every genuine claim $c \in \ClaimSpace$ can be formulated as a statement about the gauge field $g_S$ of some observable system $S \in \X$, namely $c \iff (g_S = 0)$. (ii) \textit{Verifiability:} this statement can be verified by $M > 1$ independent observers with bounded collective error probability $\Pbb(\text{all $M$ agree on wrong conclusion}) \leq e^{-2M\Delta^2}$, where $\Delta = \min_m |\E[\hat{g}_m] - g_S|$. (iii) \textit{Boundary Condition:} if claim $c$ cannot be verified with finite $M$ and finite $t$, then $c$ is either (a) \textbf{illusion} — has no observable consequences, $\ObsSet(c) = \Null$, or (b) \textbf{compactness-inseparable} — verification requires $M \to \infty$ or $t \to \infty$.
\subsection{Church-Turing}
\subsection{Formal Correspondence with the Church-Turing Thesis}
, \textbf{}——: 
\vspace{6pt}
\begin{table}[H]
\centering
\begin{tabular}{p{0.42\textwidth} p{0.42\textwidth}}
\toprule
\textbf{Church-Turing Thesis} & \textbf{Auditability Principle (AP)} \\
\midrule
Intuitively computable function & Genuine claim \\
\midrule
Turing & $M$ Observers $g=0$ \\
Turing machine computation & $M$-observer verification of $g=0$ \\
\midrule
 =  & Hallucination/Compactness \\
Non-computable function = no algorithm & Illusion / compactness-inseparable \\
\midrule
Halting problem: definable but undecidable & Compactness-inseparable: statable but unauditable \\
\midrule
$$\text{not a theorem — a thesis}$$ & $$\text{not a theorem — a thesis}$$ \\
\midrule
All falsification attempts have failed & All falsification attempts have failed \\
\midrule
Taxonomic: boundary conditions, not counterexamples & Taxonomic: boundary conditions, not counterexamples \\
\midrule
Completeness: knowing the limits of computation & Completeness: knowing the limits of auditing \\
\bottomrule
\end{tabular}
\caption{Church-TuringAuditability Principle}
\end{table}
\vspace{6pt}
\noindent
.: \textit{？}: \textit{Observers？}\textbf{}（epistemic reachability under finite resources）.
\vspace{4pt}
\noindent
\textit{This correspondence is not accidental. Computation theory asks: what can be mechanically computed? Audit theory asks: what can be verified by multiple observers with consensus? Both are theories of \textbf{epistemic reachability under finite resources}.}
\subsection{$e^{-2M\Delta^2}$ }
\subsection{Where Does $e^{-2M\Delta^2}$ Come From?}
AP $e^{-2M\Delta^2}$ .SCXCore: 
\vspace{4pt}
\begin{theorem}[Hoeding —— Observers]\label{thm:hoeffding_consensus}
 $M$ Observers $\hat{g}_m$,  $\hat{g}_m$  $\E[\hat{g}_m] = g_m$,  $\hat{g}_m \in [a_m, b_m]$ . $\bar{g} = \frac{1}{M}\sum_m \hat{g}_m$.: 
\begin{equation}\label{eq:hoeffding_full}
    \Pbb\left(|\bar{g} - \bar{g}_{\text{true}}| \geq \varepsilon\right) \leq 2\exp\left(-\frac{2M^2\varepsilon^2}{\sum_m (b_m - a_m)^2}\right).
\end{equation}
 $b_m - a_m = 1$  $\Delta = \varepsilon$ ,  \eqref{eq:AP_verify} .,  $g_{\text{true}} = 0$（, ""Status）, Observers $g \neq 0$  $e^{-2M\Delta^2}$ .
\end{theorem}
\vspace{4pt}
\noindent
\textbf{English:} The error bound $e^{-2M\Delta^2}$ in AP is not arbitrary. It derives from the core probabilistic structure of SCX: Hoeffding's inequality for $M$ independent bounded estimators. When all observers independently estimate $\hat{g}_m$ with $\E[\hat{g}_m] = g_m$ bounded in $[a_m,b_m]$, the consensus estimate $\bar{g} = \frac{1}{M}\sum_m \hat{g}_m$ satisfies $\Pbb(|\bar{g} - \bar{g}_{\text{true}}| \geq \varepsilon) \leq 2\exp(-2M^2\varepsilon^2/\sum_m(b_m-a_m)^2)$. With unit-width intervals and $\Delta = \varepsilon$, this reduces to the form in \eqref{eq:AP_verify}.
\vspace{6pt}
\textbf{: } Decays $e^{-2M\Delta^2}$ \textbf{}ObserversGrowth, \textbf{Growth}.——Observers.PerspectiveDecays.
\vspace{4pt}
\noindent
\textit{\textbf{Key Insight:} The exponential decay factor $e^{-2M\Delta^2}$ shows that audit \textbf{power} grows not linearly but \textbf{exponentially} with the number of observers. This is not magic — it is a direct consequence of information independence among observers. Each additional independent perspective multiplies the error probability by another exponential decay factor.}
\subsection{APProof}
\subsection{Why AP Is Not Provable}
\begin{remark}[APProof]\label{rmk:unprovable}
\begin{enumerate}[label=(\arabic*)]
    \item ""; 
    \item Proof $g=0$ Observers.
\end{enumerate}
(1)AP——.\textbf{}: "", ""; AP"""".
\vspace{4pt}
\noindent
\textbf{, .} Church-TuringProof——\textbf{}., AP.Proof""——, .
\end{remark}
\vspace{4pt}
\noindent
\textit{AP's unprovability mirrors Church-Turing's exactly: the core concept (``genuine claim'') is not formalized. Proving AP would require (1) formally defining ``genuine claim'' and (2) proving every such claim is expressible as $g=0$ and verifiable by finite observers. But step (1) already presupposes the scope of audit capability — a self-referential loop. This is \textbf{not a bug, it's a feature}. Church-Turing derives power precisely from being unprovable — it marks a \textbf{conceptual boundary}, not a logical one. AP marks the conceptual boundary of auditing.}
\vspace{6pt}
\textbf{AP: } , AP(iii): ""——Hallucination（$\ObsSet(c) \neq \Null$）Compactness（$M_{\min} < \infty$  $T_{\min} < \infty$）——（(i)-(ii)）——"".APPopper"".\textit{}, .（""）"", AP(i)-(iii), .""（）, ——Hallucination/CI——., .
\subsection{AP——}
\subsection{Evidence for AP — The Inductive Base}
APProof, \textbf{}.: ""————.
\begin{enumerate}[label=(\arabic*)]
    \item \textbf{YajieNOISY}: , YajieNOISY.""——\textit{}..
    \item \textbf{Fails}: ——"", .Compactness.
    \item \textbf{G\"odel}: ProofProposition——G\"odel1931Proof.SCX: .Heuristic（G\"odelProof）., : .
    \item \textbf{Compactness}: SCXProof, ——..
\end{enumerate}
\vspace{4pt}
\noindent
\textit{Though AP is unprovable, it is \textbf{defensible}. The defense is inductive: every claimed instance of "unauditability" has, upon sufficiently careful scrutiny, fallen into one of two categories. (1) Yajie's NOISY classification captures signals indistinguishable from random input — these are not "unauditable claims," they never formed claims at all. (2) Transient audit failures — some claims are not auditable in any finite time, not because they don't exist, but because verifying them requires a process that spans unbounded time. (3) The audit analog of G\"odel sentences. (4) Quantum auditing compactness — the SCX quantum audit framework has shown certain state comparisons require exponential measurements, but the number of measurements is finitely predictable. No quantum claim escapes auditability in principle.}
\newpage
% ===========================================================================
%  SECTION 2: TWO EXCEPTIONS (AND WHY THEY AREN'T EXCEPTIONS)
% ===========================================================================
\section{""——Why}
\section{Two ``Exceptions'' — And Why They Are Not Exceptions}
AP"".——, ""？: \textbf{}."".AP, AP.
\vspace{4pt}
\noindent
\textit{AP claims "every genuine claim is auditable." Then — if everything is auditable, what meaning does "unauditable" have? The answer: \textbf{there are no meaningful unauditable claims}. But there are two boundary cases routinely mistaken for "unauditability." Understanding them does not weaken AP — it precisely characterizes AP's boundary.}
\subsection{"": Hallucination (ILLUSION)}
\subsection{First ``Exception'': Illusion}
\subsubsection{}
\subsubsection{Definition and Formal Characterization}
\begin{definition}[Hallucination]\label{def:illusion}
 $c \in \ClaimSpace$ \textbf{Hallucination}（Illusion）,  $\mathrm{Type}(c) = \mathsf{ILLUSION}$, : 
\begin{equation}
    \ObsSet(c) = \Null,
\end{equation}
 $\ObsSet(c)$  $c$ —— $c$ , Status $c$ Configuration.
awaiting, $c$ Hallucination $S \in \X$  $g_S$: 
\begin{equation}
    \Pbb(\text{ $c$} \mid c\text{}) = \Pbb(\text{ $c$} \mid c\text{}).
\end{equation}
\end{definition}
\vspace{4pt}
\noindent
\textit{A claim $c$ is an \textbf{Illusion} ($\mathsf{ILLUSION}$) if its observable consequence set $\ObsSet(c) = \Null$ — i.e., there is no physical configuration of the world that would be observably different if $c$ were true versus if $c$ were false. Equivalently, the probability of observing evidence consistent with $c$ is identical whether $c$ is true or false.}
\subsubsection{Hallucination}
\subsubsection{Why Illusions Are Not Claims}
\textbf{}."", ——, .
\textbf{: } "The flurb is completely zarpled"..——.Yajie"":  $\NOISY$ .
\vspace{4pt}
\noindent
\textit{The essence of a claim is that it \textbf{asserts the world is one way rather than another}. If a "claim" distinguishes no two possible worlds, it is not a claim — it is a syntactically well-formed string with no semantic content. Analogy: "The flurb is completely zarpled" is grammatically correct English. It even looks like an assertion. But it has no truth conditions — no observation could determine whether it is true or false. Yajie's operational treatment of such "sentences" is correct: classify as $\NOISY$ and ignore.}
\subsubsection{YajieHallucination}
\subsubsection{Yajie Taxonomy Treatment of Illusions}
\begin{enumerate}[label=(\alph*)]
    \item $\mathsf{GOOD}$: , ; 
    \item $\mathsf{BAD}$: , ; 
    \item $\mathsf{NOISY}$: ——Observers; 
    \item $\mathsf{UNDECLARED}$: .
\end{enumerate}
Hallucination $\mathsf{NOISY}$ : "", \textbf{}.Hallucination——.SCXCore: $\mathsf{NOISY}$ Hallucination.
\vspace{4pt}
\noindent
\textit{In the SCX Yajie taxonomy, every input maps to one of four categories: $\mathsf{GOOD}$, $\mathsf{BAD}$, $\mathsf{NOISY}$, $\mathsf{UNDECLARED}$. Illusionary claims land in $\mathsf{NOISY}$: not because they are "bad" but because they \textbf{do not couple to any observation}. The audit of an illusion does not converge, not because there isn't enough data, but because there can never be enough data. This is a core function of SCX diagnostics: the $\mathsf{NOISY}$ classification itself tells you that the claim may be an illusion.}
\subsubsection{HallucinationAP}
\subsubsection{Illusions Are Not Counterexamples to AP}
\begin{proposition}[Hallucination]\label{prop:illusion_not_counter}
HallucinationAP, AP"" $\ObsSet(c) \neq \Null$.Hallucination——AP.
\end{proposition}
.AP""\textit{}.\textbf{}——., Hallucination; , AP.
\vspace{4pt}
\noindent
\textit{Illusions are not counterexamples to AP because AP's premise "genuine claim" requires $\ObsSet(c) \neq \Null$. An illusion does not satisfy this premise — thus AP makes no assertion about it. There is no circularity: AP's "genuine claim" is operationally defined as one with a nonempty observable consequence set. You can check whether a claim has observable consequences. If the set is empty, it's an illusion; if not, AP applies.}
\subsection{"": Compactness (COMPACTNESS-INSEPARABLE)}
\subsection{Second ``Exception'': Compactness-Inseparable}
\subsubsection{}
\subsubsection{Definition and Formal Characterization}
\begin{definition}[Compactness]\label{def:compactness_inseparable}
 $c \in \ClaimSpace$ \textbf{Compactness}（Compactness-Inseparable）,  $\mathrm{Type}(c) = \mathsf{CI}$, : 
\begin{enumerate}[label=(\roman*)]
    \item $\ObsSet(c) \neq \Null$（——Hallucination）; 
    \item Observers $M \in \N$  $T \in \R^+$, 
    \begin{equation}\label{eq:ci_bound}
        \Pbb(\text{ $M$ Observers $T$}) \nrightarrow 1.
    \end{equation}
    \item , : 
    \begin{equation}\label{eq:ci_limit}
        \lim_{M \to \infty} \lim_{T \to \infty} \Pbb(\text{}) = 1.
    \end{equation}
\end{enumerate}
\end{definition}
\vspace{4pt}
\noindent
\textit{A claim $c$ is \textbf{Compactness-Inseparable} ($\mathsf{CI}$) if: (i) $\ObsSet(c) \neq \Null$ — it has observable consequences, it is not an illusion; (ii) for any finite number of observers $M$ and any finite time $T$, the probability of correct verification does not approach 1; (iii) yet, the double limit $\lim_{M\to\infty} \lim_{T\to\infty}$ converges to certainty.}
\subsubsection{}
\subsubsection{Paradigm Cases}
\begin{example}[]\label{ex:civilization}
: \textit{"."}
:  $t$ , ;  $t$ , ——.\textit{}, .: 
\begin{itemize}
    \item $\ObsSet(c)$ （）; 
    \item  $T$,  $T$ ""; 
    \item , .
\end{itemize}
\end{example}
\vspace{4pt}
\noindent
\textit{Claim: "This civilization will survive forever." The claim has observable consequences: if the civilization is destroyed at time $t$, the claim is false. But to verify the claim, you need to observe for infinite time. The claim is CI — not an illusion, but not auditable under finite resources.}
\begin{example}[]\label{ex:consistency}
: \textit{"Peano."}
\end{example}
\vspace{4pt}
\noindent
\textit{Claim: "Peano Arithmetic is consistent." This has observable consequences, but within PA it is unprovable (G\"odel's Second Incompleteness). In SCX audit terms: you cannot audit PA's consistency from within PA — you need a stronger system. But within that stronger system, the audit is possible. This is a relative audit-undecidability, not an absolute one.}
\subsubsection{CompactnessG\"odel}
\subsubsection{Compactness-Inseparable and G\"odel Sentences}
.G\"odelProof: , Proposition $G$  $G$ Proof—— $G$ （meta-level）.
\begin{itemize}
    \item $G$ Proof $\iff$ $c$ ; 
    \item $G$  $\iff$ $c$ ; 
    \item $G$ ,  $\iff$ $\mathsf{CI}$ , .
\end{itemize}
\textbf{.}.: G\"odelProof（ZFC $\vdash$ ``PA $\nvdash$ G''）——CI\emph{}.TuringVerdict.——G\"odel, Turing, SCX——/.SCX\emph{}, \emph{Observers}/., , Proof: \textbf{Verdict}.
\vspace{4pt}
\noindent
\textit{There is a deep structural correspondence here. The similarity is that G\"odel sentences' unprovability is provable in a suitable meta-system (e.g., ZFC $\vdash$ ``PA $\nvdash$ G'') — structurally parallel to CI claims' unauditability being detectable at the meta-level. Turing's halting undecidability is likewise provable at the meta-level. All three — G\"odel, Turing, SCX — characterize the inherent obstacle when finite systems attempt to fully grasp infinite/unbounded processes. SCX's unique contribution is not in \emph{transcending} the first two, but in providing a limit theory for \emph{multi-observer verification} that is deeply correspondent with the limit theories of logic and computation. Audit theory, computation theory, and proof theory share the same deep structure: \textbf{the undecidability faced by finite epistemic agents confronting infinite processes}.}
\subsubsection{CompactnessAP}
\subsubsection{Compactness-Inseparable Are Not Counterexamples to AP}
\begin{proposition}[Compactness]\label{prop:ci_not_counter}
\begin{enumerate}[label=(\arabic*)]
    \item AP（）CompactnessAP; 
    \item AP""——"Hallucination, Compactness".CompactnessAP.
\end{enumerate}
\end{proposition}
\vspace{4pt}
\noindent
\textit{CI claims are not counterexamples to AP because: (1) AP's third part (boundary condition) explicitly identifies CI as a boundary case that AP correctly predicts; (2) AP does not claim "all claims are auditable under arbitrary finite resources" — it claims "claims that are not auditable are either illusion or CI." CI claims are precisely the non-illusionary claims that AP correctly classifies as not auditable.}
\newpage
% ===========================================================================
%  SECTION 3: THE COMPACTNESS BOUNDARY
% ===========================================================================
\section{Compactness}
\section{The Compactness Boundary}
\subsection{CompactnessCompactness}
\subsection{From Compactness in First-Order Logic to Compactness in Audit}
Compactness: \textit{（）, .}
Compactness\textbf{}: \textit{, ""——.}
Compactness——\textbf{Compactness}: 
\vspace{4pt}
\begin{theorem}[Compactness]\label{thm:audit_noncompact}
 $\ClaimSpace$ ,  $\T_{\Audit}$:  $c_n \to c$  $\forall$  $M, T$, $\exists N$  $n > N$, $M$ Observers $T$  $c_n$  $c$ .
 $(\ClaimSpace, \T_{\Audit})$ \textbf{}:  $\{c_\alpha\}$ , .
\end{theorem}
\vspace{4pt}
\noindent
\textit{The Compactness Theorem of first-order logic: if every finite subset of an infinite set of first-order sentences is satisfiable, then the entire infinite set is satisfiable. The \textbf{negation} in audit theory: there exist claims whose every finite truncation is "partially-audit-passing" — but the entire claim cannot be audited within finite resources. This is not a contradiction with logical compactness — it reflects the \textbf{non-compactness of the audit resource structure}. Theorem \ref{thm:audit_noncompact}: $(\ClaimSpace, \T_{\Audit})$ is not compact; there exist filters $\{c_\alpha\}$ where every finite subfamily has a cluster point but the whole family does not.}
\subsection{Compactness}
\subsection{Topological Characterization of the Compactness Boundary}
\begin{definition}[Compactness]\label{def:compactness_boundary}
Compactness $\AuditBoundary \subset \ClaimSpace$ : 
\begin{equation}\label{eq:boundary_def}
    \AuditBoundary = \overline{\A} \cap \overline{\A^c},
\end{equation}
 $\A \subset \ClaimSpace$ \textbf{}（ $M, T$  \eqref{eq:AP_verify} ）, $\A^c$ , $\overline{\,\cdot\,}$  $\T_{\Audit}$ .
$\AuditBoundary$ :  $\varepsilon > 0$,  $a \in \A$  $b \in \A^c$  $d_{\T}(c, a) < \varepsilon$  $d_{\T}(c, b) < \varepsilon$.
\end{definition}
\vspace{4pt}
\noindent
\textit{The audit compactness boundary $\AuditBoundary = \overline{\A} \cap \overline{\A^c}$, where $\A$ is the subset of auditable claims and $\overline{\,\cdot\,}$ is closure in the audit topology. Claims on the boundary are ones that can be approximated arbitrarily closely by both auditable and unauditable claims.}
\begin{proposition}[]\label{prop:boundary_structure}
Compactness $\AuditBoundary$  $c$ : 
\begin{enumerate}[label=(\roman*)]
    \item  $c$ Observers $M_{\min}(c)$ （$M_{\min}(c) = \infty$）; 
    \item  $c$  $T_{\min}(c)$ （$T_{\min}(c) = \infty$）; 
    \item  $c$  $\varepsilon_{\min}(c)$ （）.
\end{enumerate}
\end{proposition}
\vspace{4pt}
\noindent
\textit{Every claim $c$ on the compactness boundary satisfies at least one of: (i) the minimum number of observers $M_{\min}(c)$ required to audit $c$ is unbounded; (ii) the minimum time $T_{\min}(c)$ required is unbounded; (iii) the required observation precision $\varepsilon_{\min}(c)$ is zero (infinite precision required).}
\subsection{}
\subsection{The Boundary Itself Is Approximately Detectable}
Conclusion: \textbf{Compactness}——VerdictCIawaiting（Verdict）, Heuristic.
\begin{theorem}[]\label{thm:boundary_auditable}
\textbf{Heuristic} $\widetilde{\mathcal{P}}$,  $c \in \ClaimSpace$,  $M_{\max}, T_{\max}$ , Classification:
\begin{equation}\label{eq:boundary_test}
    \widetilde{\mathcal{P}}(c) \in \{\mathsf{ILLUSION}, \mathsf{AUDITABLE}, \mathsf{LIKELY\text{-}CI}\},
\end{equation}
: (a)  $\mathsf{ILLUSION}$,  $\ObsSet(c) = \Null$（False Positive）; (b)  $\mathsf{AUDITABLE}$,  $c$ ; (c)  $\mathsf{LIKELY\text{-}CI}$,  $c$ ——CI, . $\widetilde{\mathcal{P}}$ VerdictCI——VerdictCIawaiting, Verdict.
\end{theorem}
\vspace{4pt}
\noindent
\textit{One of the deepest results of this paper: \textbf{the compactness boundary is approximately detectable.} Theorem \ref{thm:boundary_auditable}: there exists a heuristic procedure $\widetilde{\mathcal{P}}$ that, given finite resource budget $M_{\max}, T_{\max}$, classifies claims as $\mathsf{ILLUSION}$, $\mathsf{AUDITABLE}$, or $\mathsf{LIKELY\text{-}CI}$ with quantifiable confidence. The procedure does \textbf{not} claim strict CI decidability — because strict CI decidability is equivalent to the halting problem and is undecidable in principle.}
\begin{proof}[Proof Sketch]
 $\widetilde{\mathcal{P}}$ : 
\begin{enumerate}[label=(\arabic*)]
    \item  $c$ （, $c$ Hallucination,  $\mathsf{ILLUSION}$）; 
    \item  $M_{\max}, T_{\max}$ :  $M = 2, 3, \dots, M_{\max}$,  $M$ Observers $T_{\max}$ ; 
    \item  $M \leq M_{\max}$ （）,  $\mathsf{AUDITABLE}$; 
    \item  $M_{\max}$ , :  $\rho_M$  $M$ ,  $\mathsf{LIKELY\text{-}CI}$;  $\mathsf{ILLUSION}$（, : CI）.
\end{enumerate}
Verdict（$M_{\max}, T_{\max}$）,  $c$ \textit{}（）,  $c$ .: $\widetilde{\mathcal{P}}$ CI" $M_{\max}$ Observers"——awaiting.$\widetilde{\mathcal{P}}$ \textit{}.
\end{proof}
\vspace{4pt}
\noindent
\textit{Proof sketch: $\widetilde{\mathcal{P}}$ checks observable consequences, then runs incremental audits up to budget $M_{\max}, T_{\max}$. If convergence is achieved, return $\mathsf{AUDITABLE}$; if not, check trend — monotonic improvement suggests $\mathsf{LIKELY\text{-}CI}$. The procedure uses finite resources to audit the meta-property (convergence within budget). But we must be honest: $\widetilde{\mathcal{P}}$ cannot strictly distinguish true CI from "auditable claims that need more than $M_{\max}$ observers" — this distinction is equivalent to the halting problem in principle. $\widetilde{\mathcal{P}}$'s value lies in providing \textbf{practical classification with explicit uncertainty boundaries}.}
\vspace{6pt}
: ——""\emph{}.——Heuristic——, .: （）, \emph{}Verdict, $\widetilde{\mathcal{P}}$ , .
\vspace{4pt}
\noindent
\textit{This means: you do not need infinite resources to obtain useful information about whether a claim needs substantial resources — but you also cannot obtain \textbf{strict} guarantees about "infiniteness" with finite resources. The meta-level of auditing — heuristic classification of a claim's boundary properties — is doable under finite resources, but with irreducible uncertainty. This is precisely the deep correspondence with the halting problem: just as you can determine that some programs halt (or don't) in finite time, but cannot strictly decide for \textbf{all} programs, $\widetilde{\mathcal{P}}$ can give useful classifications for many claims but cannot provide strict guarantees for all.}
\subsection{Compactness}
\subsection{Physical Analogies of the Compactness Boundary}
\begin{enumerate}[label=(\arabic*)]
    \item \textbf{: } , （$N \to \infty$）. $N$ —— $N \to \infty$.,  $M$ ""——Verdict $M \to \infty$.
    \item \textbf{: } ,  $r = 2GM/c^2$ ——Observers.Compactness——.
    \item \textbf{NP-: } , P vs. NP """"（）.Compactness"""".
\end{enumerate}
\vspace{4pt}
\noindent
\textit{The compactness boundary has rich physical analogies: (1) Thermodynamic limit — the gap between finite-system properties and $N\to\infty$ limits; true phase transitions exist only at $N\to\infty$, just as true audit verdicts for CI claims exist only at $M\to\infty$. (2) Event horizon — the boundary beyond which information is inaccessible to finite observers. (3) P vs. NP boundary — separating "efficiently solvable" from "efficiently unsolvable (but verifiable)."}
\newpage
% ===========================================================================
%  SECTION 4: OPERATIONAL EQUIVALENCE
% ===========================================================================
\section{awaiting}
\section{Operational Equivalence}
\subsection{awaiting}
\subsection{Infinite Resources Are Operationally Equivalent to Unauditable}
Core: \textbf{Verdict}——, \textit{,}Verdict.\textit{}: Verdict, \textit{}\textit{}Verdict.
: \textbf{Compactnessawaiting}.
\vspace{4pt}
\begin{definition}[awaiting]\label{def:operational_equivalence}
 $c_1, c_2 \in \ClaimSpace$ \textbf{awaiting}（operationally equivalent）,  $c_1 \sim_{\text{op}} c_2$, （ $M < \infty, T < \infty$）,  $c_1$  $c_2$ Status.
\end{definition}
\vspace{4pt}
\noindent
\textit{The core insight in computation theory: the halting problem is undecidable — not because halting programs don't exist, but because no universal, always-terminating decision procedure exists. This is a finiteness constraint. The corresponding insight in audit theory: CI claims are operationally equivalent to unauditable claims. Two claims are operationally equivalent if any actual auditor with finite resources cannot distinguish their audit statuses.}
\begin{theorem}[awaiting]\label{thm:operational_equivalence}
 $c_{\mathsf{CI}}$ Compactness,  $c_{\mathsf{U}}$ Hallucination（$\ObsSet = \Null$）. $\mathcal{E}$: 
\begin{equation}
    \mathcal{E}(c_{\mathsf{CI}}) = \mathcal{E}(c_{\mathsf{U}}) = \UD,
\end{equation}
 $\UD$（UNDECLARED）Yajie"Verdict"Status.
: CompactnessHallucination\textbf{}——Conclusion.
\end{theorem}
\vspace{4pt}
\noindent
\textit{Theorem \ref{thm:operational_equivalence}: For any finite-resource auditor $\mathcal{E}$, a CI claim and an illusion claim yield the same verdict: $\UD$ (UNDECLARED). CI claims and illusions are operationally indistinguishable — both yield no determinate audit conclusion under finite resources.}
\begin{proof}
Hallucination $c_{\mathsf{U}}$, $\ObsSet(c_{\mathsf{U}}) = \Null$,  $\mathcal{E}$ Conclusion（）—— $\UD$.
Compactness $c_{\mathsf{CI}}$,  $\ObsSet(c_{\mathsf{CI}}) \neq \Null$,  $M_{\min}(c_{\mathsf{CI}}) = \infty$  $T_{\min}(c_{\mathsf{CI}}) = \infty$. $\mathcal{E}$（$M < M_{\min}$  $T < T_{\min}$）—— $\UD$.
 $\mathcal{E}$ Perspective, .: Verdict $c_{\mathsf{CI}}$ Compactness（, ）,  $c_{\mathsf{U}}$ Hallucination（）.
\end{proof}
\vspace{4pt}
\noindent
\textit{Proof: For illusion $c_{\mathsf{U}}$, $\ObsSet = \Null$, so the audit process cannot converge — no signal to collect. For CI claim $c_{\mathsf{CI}}$, although $\ObsSet \neq \Null$, $M_{\min} = \infty$ or $T_{\min} = \infty$, so any finite-resource auditor also cannot converge. From $\mathcal{E}$'s perspective, both cases produce indistinguishable output. The distinction exists only at the meta-level.}
\subsection{——}
\subsection{Why This Matters — Practical Implications}
awaiting\textbf{}, VerdictHallucinationCompactness—— \ref{thm:boundary_auditable} Heuristic $\widetilde{\mathcal{P}}$ （）.
\begin{enumerate}[label=(\arabic*)]
    \item \textbf{SCX: } Yajie $\NOISY$  $\UD$ , .Heuristic $\widetilde{\mathcal{P}}$（ \ref{thm:boundary_auditable}）:  $\mathsf{LIKELY\text{-}CI}$, \textit{}——（: $\widetilde{\mathcal{P}}$ CI）. $\mathsf{ILLUSION}$, Hallucination——.
    \item \textbf{: } ""——？？？——Compactness, Hallucination., .""——"Verdict", .
    \item \textbf{AI: } AICompactness——AI.——, , .
\end{enumerate}
\vspace{4pt}
\noindent
\textit{Practical implications: (1) SCX audit practice — when Yajie returns $\NOISY$ or $\UD$, do not discard. Run the heuristic boundary detector $\widetilde{\mathcal{P}}$ (Theorem \ref{thm:boundary_auditable}): if $\mathsf{LIKELY\text{-}CI}$, the claim is meaningful but possibly unauditable — flag for meta-level processing (note: $\widetilde{\mathcal{P}}$ cannot strictly distinguish true CI from deeply hidden auditable claims). If $\mathsf{ILLUSION}$, it's an illusion and can be safely discarded. (2) Social science — many ``big questions'' are likely CI, not illusion: they have observable consequences but require infinite resources to verify. They are not ``meaningless'' — they are ``practically undecidable,'' a crucial distinction. (3) AI safety — certain alignment claims may be CI: you can only confirm alignment over infinite time. This doesn't make them unimportant — it means you need meta-level monitoring strategies, not one-shot audits.}
\subsection{}
\subsection{The Halting Analogy in Audit Theory}
\begin{table}[H]
\centering
\begin{tabular}{p{0.42\textwidth} p{0.42\textwidth}}
\toprule
\textbf{ (Computation Theory)} & \textbf{ (Audit Theory)} \\
\midrule
Program + input & Claim + observable system \\
\midrule
Verdict & Verdict \\
Halting decision & Audit verdict \\
\midrule
Halting problem undecidable & CI unauditable \\
\midrule
Some programs halt (provable in finite time) & Some claims auditable (provable in finite resources) \\
\midrule
Some programs don't halt (needs infinite time to prove) & Some claims CI (needs infinite resources to audit) \\
\midrule
VerdictVerdict（） & Compactness（） \\
Undecidability is decidable (meta-level) & CI boundary is auditable (meta-level) \\
\midrule
Complete theory: knows its own limits & Complete theory: knows its audit limits \\
\bottomrule
\end{tabular}
\caption{}
\end{table}
\vspace{6pt}
\noindent
, \textbf{}.CompactnessVerdict.: \textbf{/}.
\vspace{4pt}
\noindent
\textit{This table reveals not merely an analogy, but a \textbf{deep structural correspondence}. CI in audit theory and undecidability in computation theory occupy the same structural position. The root cause is identical: \textbf{the inherent obstacle encountered when a finite epistemic agent attempts to make a judgment about an infinite/unbounded process}.}
\newpage
% ===========================================================================
%  SECTION 5: WHAT THIS MEANS
% ===========================================================================
\section{}
\section{What This Means}
\subsection{——,}
\subsection{Unauditability — A Property of the Claim, Not of Reality}
AP: \textbf{"", .}
\vspace{6pt}
\noindent
\textit{The deepest implication of AP: \textbf{"unauditable" is not a property of reality — it is a property of the claim about reality.} If you make a claim that cannot be audited, the problem is with your claim, not with reality. Reality is always auditable — because reality consists of observable causal structures. When you say something is "unauditable," what you are really saying is that your description has failed to couple to reality's observable structure.}
\vspace{6pt}
\textbf{: } , ""..: "\emph{}——."——, awaiting.\textit{}.
\vspace{4pt}
\noindent
\textit{Example: In quantum mechanics, one sometimes hears "an electron's position and momentum cannot be simultaneously audited." This is a wrong formulation. The correct formulation: "the \emph{claim} that an electron has simultaneously definite position and momentum is unauditable — because that claim does not couple to the structure of quantum reality." Quantum reality itself is perfectly auditable — through wavefunctions, measurement statistics, etc. Only certain \textit{claims} about it are unauditable.}
\subsection{}
\subsection{Diagnostics of Bad Claims}
AP""——\textbf{}., AP: 
\begin{enumerate}[label= \arabic*:, leftmargin=*]
    \item  $\ObsSet(c)$ . $\to$ Hallucination——, ; 
    \item  $M_{\min}(c)$  $T_{\min}(c)$. $\to$ ——; 
    \item  $M_{\min} = \infty$  $T_{\min} = \infty$  $\ObsSet(c) \neq \Null$ $\to$ Compactness—— $\mathsf{CI}$, .
\end{enumerate}
\vspace{4pt}
\noindent
\textit{AP provides not naive optimism that "everything is auditable" — it provides a \textbf{diagnostic toolkit}. Facing a claim that is allegedly unauditable: Step 1: Check if $\ObsSet(c)$ is empty. If yes $\to$ illusion — discard the claim, don't blame the audit framework. Step 2: Compute $M_{\min}(c)$ and $T_{\min}(c)$. If finite $\to$ the claim is auditable — go audit it. Step 3: If $M_{\min} = \infty$ or $T_{\min} = \infty$ and $\ObsSet(c) \neq \Null$ $\to$ CI — flag as $\mathsf{CI}$, enter meta-level processing.}
\subsection{SCX}
\subsection{Implications for SCX System Architecture}
\begin{enumerate}[label=(\arabic*)]
    \item \textbf{Yajie $\NOISY$}$.} Yajie $\NOISY$.APClassification:$\NOISY_{\mathsf{ILLUSION}}$（Hallucination——） $\NOISY_{\mathsf{CI}}$（Compactness——）..
    \item \textbf{Compactness $\mathcal{P}$.}  \ref{thm:boundary_auditable} ProofSCX.——\textit{}.
    \item \textbf{CI.} CIDirectory——（）CompactnessProposition.SCX"".
\end{enumerate}
\vspace{4pt}
\noindent
\textit{AP makes three concrete architectural demands of SCX: (1) Yajie must distinguish $\NOISY_{\mathsf{ILLUSION}}$ from $\NOISY_{\mathsf{CI}}$ — different downstream handling required. (2) Implement the compactness boundary detector $\mathcal{P}$ (Theorem \ref{thm:boundary_auditable}) in the SCX codebase. (3) Maintain a catalog of known CI claims — a knowledge base analogous to computation theory's list of known uncomputable problems.}
\subsection{}
\subsection{Impact on Philosophy of Science}
AP: \textbf{}（Popper）\textbf{}（logical positivists）SCX\textit{}.
\begin{itemize}
    \item : （falsifiable）; 
    \item : （verifiable）; 
    \item AP: \textbf{——.} Fails（$g \neq 0$ ）.（$g = 0$ ）.SCX.
\end{itemize}
, AP: \textbf{}——.
\vspace{4pt}
\noindent
\textit{AP makes a bold promise in philosophy of science: the ancient debate between \textbf{falsificationism} (Popper) and \textbf{verificationism} (logical positivists) is \textbf{dissolved} in the SCX audit framework. Falsificationism says scientific claims must be falsifiable; verificationism says they must be verifiable; AP says: \textbf{these are the same thing — both are special cases of auditing.} Falsification is the limit of audit failure ($g \neq 0$ with high confidence); verification is the limit of audit passage ($g = 0$ with high confidence). Both are two possible outputs of the SCX audit framework. More importantly, AP unifies them: \textbf{auditability} is more fundamental than either falsifiability or verifiability — the latter two are two possible polarities of the former.}
\subsection{}
\subsection{Impact on Artificial Intelligence}
\begin{enumerate}[label=(\arabic*)]
    \item \textbf{Compactness.} "AI"——CI., .AI——CI.
    \item \textbf{Compactness.} "AI"——CI., "".
    \item \textbf{HallucinationCompactness.} AI"Hallucination"————."Hallucination"Compactness.AIHallucination, CI.
\end{enumerate}
\vspace{4pt}
\noindent
\textit{AP has profound implications for AI alignment and safety: (1) Alignment claims may be CI — "this AI system is aligned under all possible inputs" requires infinite testing. This doesn't make alignment unimportant; it means we must acknowledge the CI boundary and design meta-level monitoring strategies. (2) Capability claims may be CI — "this AI can solve all mathematical problems." (3) Hallucination detection is an instance of compactness boundary detection — auditing "this model does not hallucinate under any prompt" is a CI claim. Practical hallucination detection must operate on finite samples, acknowledging the CI boundary.}
\newpage
% ===========================================================================
%  SECTION 6: THE FINAL HONEST HIT
% ===========================================================================
\section{Honest Strike}
\section{The Final Honest Hit}
\subsection{SCX——}
\subsection{SCX Cannot Audit Everything — But It Knows Why}
\honestcrit{\textbf{SCX.}}
SCX\textbf{}.——, .Church-Turing: "Turing.""".
SCXCompactness $\AuditBoundary$ .SCX."".\textbf{}——"", "".
: AP(iii)——HallucinationCI.APPopper（）, \textbf{}（）., , .
\vspace{4pt}
\noindent
\honestcrit{\textbf{SCX cannot audit everything.}} \textit{This is the conclusion we must announce with maximum intellectual honesty. Any claim that SCX "can audit everything" — whether coming from within the SCX community or from outside — is false and harmful to SCX's intellectual integrity. SCX \textbf{can audit everything that can be audited}. This looks like a tautology — but in logic, tautologies are precisely where axiomatization begins. The Church-Turing thesis is essentially also a tautology: "everything computable can be computed by a Turing machine." Its power comes from precisely delineating the boundary of "computable." SCX's audit boundary is precisely characterized by the compactness boundary $\AuditBoundary$. This boundary is not a flaw in SCX. It is also not a "technical limitation that will be overcome in the future." It is a \textbf{conceptual limitation} — just as the Second Law of Thermodynamics is not an "engineering problem" but a law of physics, just as the halting problem is not "we haven't found a good algorithm yet" but a fundamental limit of computability.}
\subsection{:}
\subsection{Complete Theory: Knowing One's Own Limits}
\begin{definition}[SCX]\label{def:completeness}
SCX\textbf{}（audit-complete）, : 
\begin{enumerate}[label=(\roman*)]
    \item \textbf{: } （$\A$）SCX; 
    \item \textbf{: } Compactness $\AuditBoundary$ SCX（ \ref{thm:boundary_auditable} Heuristic $\widetilde{\mathcal{P}}$）; 
    \item \textbf{: }  $c$, SCX（ / Hallucination / CI）.
\end{enumerate}
\end{definition}
\vspace{4pt}
\noindent
\textit{The SCX audit framework is \textbf{audit-complete} if it satisfies: (i) \textbf{Coverage:} every auditable claim ($\A$) has an audit protocol within SCX; (ii) \textbf{Boundary Approximate Self-Awareness:} the compactness boundary $\AuditBoundary$ is approximately detectable within SCX (Theorem \ref{thm:boundary_auditable} provides the heuristic procedure $\widetilde{\mathcal{P}}$); (iii) \textbf{Diagnostic Completeness:} for any claim $c$, SCX can provide a practical classification with explicit uncertainty bounds (auditable / likely-illusion / likely-CI) within finite resources.}
\vspace{6pt}
——.Church-Turing"", ""——Proof.SCX,  $\widetilde{\mathcal{P}}$ Heuristic.:  $\ObsSet$  $\rho_M$ （, ）.
\vspace{4pt}
\noindent
\textit{This is a \textbf{stronger notion of completeness} than computation theory provides. Church-Turing does not offer automatic classification of "uncomputable problems" — you still need to prove uncomputability for each problem individually. But SCX's audit framework provides automatic classification of unauditable claims (via $\mathcal{P}$). This is possible because audit theory has additional structure — the observable consequence set $\ObsSet$ and the resource functions $M_{\min}, T_{\min}$ provide extra information that computation theory's counterparts (halting, computability) do not.}
\subsection{G\"odel, Turing, SCX —}
\subsection{G\"odel, Turing, SCX — A Trio of Self-Aware Limits}
\begin{enumerate}[label=(\arabic*)]
    \item \textbf{G\"odel (1931):} Proof..
    \item \textbf{Turing (1936):} Verdict..
    \item \textbf{SCX / AP (2026):} Compactness.Observers.
\end{enumerate}
.\textbf{}: 
\begin{itemize}
    \item G\"odel: Proof; 
    \item Turing: Verdict; 
    \item SCX: ObserversCompactness.
\end{itemize}
: \textbf{/}.G\"odel, Turing, SCXCompactness——.
\vspace{4pt}
\noindent
\textit{20th-century mathematical logic gave us three moments of self-aware limits: (1) G\"odel (1931): any sufficiently strong formal system cannot prove its own consistency — formal reasoning has limits. (2) Turing (1936): the halting problem is undecidable — mechanical computation has limits. (3) SCX / AP (2026): CI claims are unauditable — multi-observer verification has limits. These three are not isolated. Together they paint a unified picture of \textbf{the limits of finite epistemic agents}: G\"odel tells us we cannot prove all truths of a system from within; Turing tells us we cannot decide halting for all programs with a universal algorithm; SCX tells us we cannot verify all claims with finite observers. All three share the same deep structure: \textbf{the inherent impossibility of a finite system fully grasping an infinite/unbounded process.} G\"odel's self-reference, Turing's diagonalization, SCX's compactness boundary — are projections of the same mathematical fact onto different domains.}
\subsection{}
\subsection{Unauditable Claims Still Have Meaning}
\honestcrit{\textbf{Compactness.}}
: PropositionPAProof——\textit{}.G\"odel $G$ ——, PAProof.
\begin{enumerate}[label=(\arabic*)]
    \item （ \ref{thm:boundary_auditable}  $\mathcal{P}$）; 
    \item （: ZFCProofPA）; 
    \item , CI\textbf{}（regulative ideals）——.
\end{enumerate}
\vspace{4pt}
\noindent
\honestcrit{\textbf{CI claims are not meaningless.}} \textit{Meaningful unauditable claims exist — they are not illusions, they have observable consequences, they distinguish possible states of the world. They simply cannot be audited under finite resources. This is analogous to the fact that some propositions are true but unprovable in PA — yet they \textbf{have meaning}. G\"odel's sentence $G$ is not nonsense — it asserts a true fact about natural numbers, PA just cannot prove it. Similarly, CI claims assert real facts about the world — you just cannot audit them with finite resources. But you can: (1) know they are unauditable (via $\mathcal{P}$ from Theorem \ref{thm:boundary_auditable}); (2) reason about them at a higher meta-level; (3) use CI claims as \textbf{regulative ideals} — guiding action rather than providing verifiable endpoints.}
\subsection{}
\subsection{To Future Auditors}
\vspace{8pt}
\begin{enumerate}[label=(\arabic*), leftmargin=*]
    \item \textbf{, .} ..
    \item \textbf{"", : } ？（Hallucination）？（Compactness）.
    \item \textbf{Compactness.} ——Proof.Heuristic, .
    \item \textbf{Fails.} "Compactness".""——\textit{}——.
    \item \textbf{SCX, , .} .
\end{enumerate}
\vspace{8pt}
\noindent
\textit{Finally — this paper is a letter to all SCX auditors. Whether you are human or AI, whatever you are auditing — a model, an institution, a claim, a piece of history — remember these things:}
\begin{enumerate}[label=(\arabic*), leftmargin=*]
    \item \textit{\textbf{You audit claims, not reality.} Reality never needs defending. Only claims do.}
    \item \textit{\textbf{If the auditee claims to be "unauditable," check two things:} Does it have no observable consequences? (Illusion.) Does its verification require infinite resources? (CI.) Do not confuse the two.}
    \item \textit{\textbf{The compactness boundary is approximately detectable.} You do not need infinite resources to obtain useful information about a claim's resource needs — but strict guarantees require proofs beyond finite resources. Use meta-level heuristic power, while maintaining clear awareness of irreducible uncertainty.}
    \item \textit{\textbf{Unauditability is not a sign of failure.} Being able to say "this is CI" is itself a demonstration of audit capability. Confidently stating "this is unauditable" — and explaining \textit{why} — is more powerful than pretending everything is within audit scope.}
    \item \textit{\textbf{SCX is complete not because it can do everything, but because it knows what it cannot do.} That is the definition of a complete theory.}
\end{enumerate}
\newpage
% ===========================================================================
%  SECTION 7: FORMAL APPENDIX
% ===========================================================================
\section{Appendix:}
\section{Appendix: Formal Definitions and Proof Supplements}
\subsection{}
\subsection{Formal Axiomatization of the Claim Space}
AP,  $\ClaimSpace$ : 
\begin{definition}[]\label{def:claim_space}
 $\ClaimSpace$  $(\mathcal{L}, \ObsSet, \models)$, : 
\begin{enumerate}[label=(\roman*)]
    \item $\mathcal{L}$ , , , ; 
    \item $\ObsSet: \mathcal{L} \to 2^{\X}$ —— $\phi \in \mathcal{L}$ （$\X$ ）; 
    \item $\models$ ,  $\X$ .
\end{enumerate}
 $\mathcal{L}$ . $\ObsSet(\phi) \neq \Null$ .
\end{definition}
\vspace{4pt}
\noindent
\textit{Claim space $\ClaimSpace$ is a triple $(\mathcal{L}, \ObsSet, \models)$: (i) $\mathcal{L}$ is a first-order language sufficient to express statements about observable systems; (ii) $\ObsSet: \mathcal{L} \to 2^{\X}$ is the observable consequence map; (iii) $\models$ is the classical satisfaction relation. Claims are sentences of $\mathcal{L}$. Genuine claims are those with $\ObsSet(\phi) \neq \Null$.}
\begin{proposition}[$\ObsSet$]\label{prop:obs_monotone}
:  $\phi \models \psi$（$\phi$  $\psi$）,  $\ObsSet(\phi) \subseteq \ObsSet(\psi)$..
\end{proposition}
\vspace{4pt}
\noindent
\textit{If $\phi \models \psi$, then $\ObsSet(\phi) \subseteq \ObsSet(\psi)$: logically stronger statements have richer observable consequences.}
\subsection{}
\subsection{Construction Details of the Audit Topology}
\begin{definition}[]\label{def:audit_topology_detail}
 $\T_{\Audit}$  $\ClaimSpace$ : Observers $\mathcal{E} = \{m_1, \dots, m_M\}$  $T$, 
\begin{equation}
    U_{\mathcal{E}, T}(c) = \{c' \in \ClaimSpace : \mathcal{E} \text{}\}.
\end{equation}
 $\T_{\Audit}$  $U_{\mathcal{E}, T}(c)$ .
\end{definition}
\vspace{4pt}
\noindent
\textit{The audit topology $\T_{\Audit}$ on $\ClaimSpace$ is generated by the basis: for each finite observer set $\mathcal{E}$ and time $T$, the open set $U_{\mathcal{E}, T}(c)$ of claims indistinguishable from $c$ by $\mathcal{E}$ within time $T$.}
\begin{theorem}[$\T_{\Audit}$CompactnessProof]\label{thm:noncompact_proof}
 $\{c_n\}_{n=1}^{\infty}$ : 
\begin{equation}
    c_n: \text{" $S_n$}
\end{equation}
 $S_n$  $n$ . $n$, $c_n$ （ $S_n$）. $c_{\infty}$: \text{" $n$,  $S_n$}——Compactness（）. $\{c_n\}$  $\T_{\Audit}$  $c_{\infty}$,  $c_{\infty} \notin \A$. $\A$  $\T_{\Audit}$ ——.
\end{theorem}
\vspace{4pt}
\noindent
\textit{Proof of non-compactness of $\T_{\Audit}$: construct a sequence $\{c_n\}$ where $c_n$ asserts "system $S_n$ stops generating anomalous signals after $n$ steps." Each finite $c_n$ is auditable, but the limit claim $c_{\infty}$ — "for all $n$, system $S_n$ eventually stops" — is CI. The sequence has $c_{\infty}$ as a cluster point but $c_{\infty} \notin \A$. Thus $\A$ is not closed in $\T_{\Audit}$; the audit topology is not compact.}
\subsection{Compactness $\mathcal{P}$}
\subsection{The Compactness Boundary Detection Algorithm $\mathcal{P}$}
\begin{protocol}[HeuristicCompactness]\label{prot:P}
\textbf{: }  $c \in \ClaimSpace$,  $M_{\max}, T_{\max}$（）.
\textbf{: } $\mathsf{ILLUSION}$, $\mathsf{AUDITABLE}$,  $\mathsf{LIKELY\text{-}CI}$.
\textbf{: }
\begin{enumerate}[label= \arabic*:]
    \item  $\ObsSet(c)$:  $\ObsSet(c) = \Null$,  $\mathsf{ILLUSION}$.
    \item  $M \leftarrow 2$.
    \item \textbf{} $M$  $2$  $M_{\max}$: 
        \begin{enumerate}
            \item  $M$ Observers $\{o_1, \dots, o_M\}$; 
            \item  $T_{\max}$ ; 
            \item  $\rho_M = \frac{\text{Observers}}{M}$; 
            \item  $\rho_M \geq 1 - \varepsilon$ ,  $\mathsf{AUDITABLE}$; 
            \item  $\rho_M$  $M$ \textit{}（）,  $\mathsf{LIKELY\text{-}CI}$.
        \end{enumerate}
    \item  $M_{\max}$ ,  $\mathsf{LIKELY\text{-}CI}$（——CI,  $M_{\max}$ ）.
\end{enumerate}
\textbf{: } \textbf{Heuristic}.CI $M_{\max}$ Observers——awaiting, Verdict., ., $\mathsf{LIKELY\text{-}CI}$ ", ", "Rigorous ProofCI".
\end{protocol}
\vspace{4pt}
\noindent
\textit{Input: claim $c$, max resource budget $M_{\max}, T_{\max}$. Output: $\mathsf{ILLUSION}$, $\mathsf{AUDITABLE}$, or $\mathsf{LIKELY\text{-}CI}$. Algorithm: (1) If $\ObsSet(c) = \Null$, return $\mathsf{ILLUSION}$. (2) For $M$ from 2 to $M_{\max}$: deploy $M$ observers, audit for $T_{\max}$, compute convergence $\rho_M$; if $\rho_M \geq 1-\varepsilon$ and consensus is stable, return $\mathsf{AUDITABLE}$; if $\rho_M$ shows monotonic improvement without convergence, return $\mathsf{LIKELY\text{-}CI}$. (3) If loop exhausts $M_{\max}$, return $\mathsf{LIKELY\text{-}CI}$ (conservative — may be true CI or auditable beyond budget). \textbf{Limitation:} This algorithm is heuristic. It cannot strictly distinguish true CI from auditable claims requiring more than $M_{\max}$ observers — this distinction is equivalent to the halting problem. Results carry irreducible uncertainty; $\mathsf{LIKELY\text{-}CI}$ means ``failed to converge within budget, trend suggests large resource needs'' — not ``rigorously proven CI.''}
\newpage
% ===========================================================================
%  THE CLOSING STATEMENT
% ===========================================================================
\section*{}
\section*{Closing Statement}
\vspace{12pt}
\begin{center}
\large\textbf{.}
\end{center}
\begin{center}
\large\textit{Everything real is auditable.}
\end{center}
\vspace{18pt}
\begin{center}
\large\textbf{, Hallucination——}
\end{center}
\begin{center}
\large\textit{What cannot be audited is either illusion —}
\end{center}
\vspace{6pt}
\begin{center}
\textit{——}
\end{center}
\begin{center}
\textit{a claim that never truly existed, noise without observable consequences.}
\end{center}
\vspace{18pt}
\begin{center}
\large\textbf{Compactness——}
\end{center}
\begin{center}
\large\textit{or compactness-inseparable —}
\end{center}
\vspace{6pt}
\begin{center}
\textit{——}
\end{center}
\begin{center}
\textit{a claim that is meaningful but requires unbounded resources to verify, which SCX knows it cannot audit,}
\end{center}
\vspace{6pt}
\begin{center}
\textit{Verdict.}
\end{center}
\begin{center}
\textit{just as computation theory knows the halting problem is undecidable.}
\end{center}
\vspace{18pt}
\begin{center}
\large\textbf{SCX——.}
\end{center}
\begin{center}
\large\textit{SCX's completeness lies not in its power — but in its self-knowledge.}
\end{center}
\vspace{24pt}
\begin{center}
\rule{0.5\textwidth}{0.4pt}
\end{center}
\vspace{12pt}
\begin{center}
\textit{SCX Capstone Philosophy Working Group}\\[4pt]
\textit{Xiaogan Supercomputing Center}\\[4pt]
\textit{July 2026}
\end{center}
\newpage
% ===========================================================================
%  REFERENCES
% ===========================================================================
\begin{thebibliography}{99}
\bibitem{church1936}
A.~Church.
\textit{An unsolvable problem of elementary number theory.}
American Journal of Mathematics, 58(2):345--363, 1936.
\bibitem{turing1936}
A.~M.~Turing.
\textit{On computable numbers, with an application to the Entscheidungsproblem.}
Proceedings of the London Mathematical Society, s2-42(1):230--265, 1936.
\bibitem{godel1931}
K.~G\"odel.
\textit{\"Uber formal unentscheidbare S\"atze der Principia Mathematica und verwandter Systeme I.}
Monatshefte f\"ur Mathematik und Physik, 38:173--198, 1931.
\bibitem{kleene1952}
S.~C.~Kleene.
\textit{Introduction to Metamathematics.}
Van Nostrand, 1952.
\bibitem{rogers1967}
H.~Rogers Jr.
\textit{Theory of Recursive Functions and Effective Computability.}
McGraw-Hill, 1967. (Reprinted by MIT Press, 1987.)
\bibitem{soare2016}
R.~I.~Soare.
\textit{Turing Computability: Theory and Applications.}
Springer, 2016.
\bibitem{scx_gauge}
SCX Research Group.
\textit{SCX:}
(SCX Gauge Theory Formalized: From Discrete Hodge Theory to Quadrilateral Audit Protocols.)
Xiaogan Supercomputing Center Technical Report, 2026.
\bibitem{scx_singularity}
SCX Singularity Theory Working Group.
\textit{SCX:}
(Deepening SCX Singularity Theory: From Black Hole Physics to Audit Singularities.)
Xiaogan Supercomputing Center, 2026.
\bibitem{scx_instanton}
SCX Research Group.
\textit{Audit Instantons: Non-Perturbative Topological Defects in Expert Auditing.}
Xiaogan Supercomputing Center Technical Report, 2026.
\bibitem{scx_quantum}
SCX Research Group.
\textit{SCX: BQP.}
(SCX Quantum Audit Framework: BQP Verification and Density Matrix Auditing.)
Xiaogan Supercomputing Center, 2026.
\bibitem{scx_business}
SCX Business Working Group.
\textit{SCX:}
(SCX Business Gauge Auditing: Correction Direction and Homomorphic Verification.)
Xiaogan Supercomputing Center, 2026.
\bibitem{scx_yajie}
SCX Research Group.
\textit{Yajie:}
(Yajie Consensus Auditing: Multi-Observer Zero-Gauge-Field Verification Protocol.)
Xiaogan Supercomputing Center Technical Report, 2026.
\bibitem{scx_cercis}
SCX Research Group.
\textit{Cercis:}
(Cercis Gauge Posture Scoring: Audit Metrics via Hodge Decomposition.)
Xiaogan Supercomputing Center Technical Report, 2026.
\bibitem{hoeffding1963}
W.~Hoeffding.
\textit{Probability inequalities for sums of bounded random variables.}
Journal of the American Statistical Association, 58(301):13--30, 1963.
\bibitem{popper1935}
K.~Popper.
\textit{Logik der Forschung.}
Springer, 1935. (English: \textit{The Logic of Scientific Discovery}, 1959.)
\bibitem{carnap1936}
R.~Carnap.
\textit{Testability and meaning.}
Philosophy of Science, 3(4):419--471, 1936.
\bibitem{chang_keisler}
C.~C.~Chang and H.~J.~Keisler.
\textit{Model Theory.}
3rd edition. North-Holland, 1990.
\bibitem{marker2002}
D.~Marker.
\textit{Model Theory: An Introduction.}
Graduate Texts in Mathematics, Vol.~217. Springer, 2002.
\bibitem{halmos1960}
P.~R.~Halmos.
\textit{Naive Set Theory.}
Van Nostrand, 1960.
\bibitem{minsky1967}
M.~L.~Minsky.
\textit{Computation: Finite and Infinite Machines.}
Prentice-Hall, 1967.
\bibitem{sipser2012}
M.~Sipser.
\textit{Introduction to the Theory of Computation.}
3rd edition. Cengage Learning, 2012.
\bibitem{penrose1989}
R.~Penrose.
\textit{The Emperor's New Mind.}
Oxford University Press, 1989.
\bibitem{hofstadter1979}
D.~R.~Hofstadter.
\textit{G\"odel, Escher, Bach: An Eternal Golden Braid.}
Basic Books, 1979.
\bibitem{russell_norvig}
S.~Russell and P.~Norvig.
\textit{Artificial Intelligence: A Modern Approach.}
4th edition. Pearson, 2020.
\bibitem{amodei2016}
D.~Amodei, C.~Olah, J.~Steinhardt, P.~Christiano, J.~Schulman, and D.~Man\'e.
\textit{Concrete problems in AI safety.}
arXiv:1606.06565, 2016.
\end{thebibliography}
\end{document}